\hypertarget{2-related-work}{%
\section{Related Work}\label{2-related-work}}

\hypertarget{21-scres-measurement-and-adaptive-capability}{%
\subsection{SCRES measurement and adaptive capability}\label{21-scres-measurement-and-adaptive-capability}}

Supply chain resilience (SCRES) is commonly defined around the ability of
a supply chain to prepare for, respond to, and recover from disruptions
\cite{ponomarov2009understanding}. More recent theory distinguishes
resilience as a system property from resilience as a dynamic capability
that must be renewed through experience and organizational learning
\cite{wieland2021two,levitt1988organizational}. Quantitative SCRES
research has consequently moved beyond binary disruption survival toward
metrics for recovery speed, service continuity, tail performance, and
system viability \cite{hosseini2019review}. This paper follows that
measurement tradition: reinforcement learning is evaluated on
Garrido/Excel ReT and order-level recovery metrics, not on the training
reward alone.

\hypertarget{22-garrido-grounded-des-for-resilience}{%
\subsection{Garrido-grounded DES for resilience}\label{22-garrido-grounded-des-for-resilience}}

Garrido-Rios~\cite{garrido2017} provides an unusually useful starting
point for AI-enabled SCRES because the military food supply chain is
already modeled as a detailed 13-operation simulation with resilience
indices, risk families, and static policy comparisons. Garrido et
al.~\cite{garrido2024ijpr} extend that lineage through a
Cobb--Douglas-style factory resilience index, while Garrido et
al.~\cite{garrido2024iccl} explicitly identify learning algorithms as a
future path for simulation-based SCRES models. The present study answers
a narrower and more testable version of that agenda: not whether AI can
be attached to a DES, but whether learned closed-loop control can beat
strong static resilience baselines when the controllable decision surface
does or does not reach the operational recovery bottleneck.

\hypertarget{23-rl-for-supply-chain-operational-control}{%
\subsection{RL for supply-chain operational control}\label{23-rl-for-supply-chain-operational-control}}

Reinforcement learning has become an established paradigm for adaptive
decision-making in supply chain management. Yan et al.~\cite{yan2022}
and Rolf et al.~\cite{rolf2023} review applications across inventory
ordering, vehicle routing, warehouse operations, and multi-echelon
coordination. Boute et al.~\cite{boute2022} provide a design roadmap for
deep RL in inventory control, and Powell~\cite{powell2022} cautions that
neural value-function approximation is only one policy class among
several structured alternatives. Empirical inventory-control studies
reinforce that warning: deep RL can perform well, but strong base-stock,
teacher, and heuristic policies remain difficult comparators
\cite{gijsbrechts2022,demoor2022,dehaybe2024,geevers2024,temizoz2025,vanvuchelen2025}.
Our contribution is therefore not the generic application of PPO to a
supply chain. It is a controlled benchmark showing that action-surface
design can decide whether any learned policy has a useful recovery
frontier to exploit.

\hypertarget{24-des-ml-and-ai-scres}{%
\subsection{DES+ML and AI-SCRES workflows}\label{24-des-ml-and-ai-scres}}

Hybrid simulation and machine-learning workflows are now common in
supply-chain analysis. Kogler and Maxera~\cite{kogler2026} review 72
papers and 43 related studies combining discrete simulation with machine
learning, while Perez-Pena et al.~\cite{perezpena2025} show that many
simulation-ML workflows emphasize prediction or classification rather
than optimization. In resilience-oriented AI, Ding et al.~\cite{ding2026}
use MAPPO for strategic topology reconfiguration under disruption,
Bussieweke et al.~\cite{bussieweke2024recovery} combine system dynamics
and RL for recovery policies, Kim~\cite{kim2024} studies cooperative
MARL for inventory transshipment, Liu et al.~\cite{liu2025pom} apply
HAPPO to decentralized multi-echelon inventory, and Kotecha and del Rio
Chanona~\cite{kotecha2025} combine graph neural networks with MAPPO.
Other AI-SCRES studies focus on prediction, classification, or
explanation rather than online control
\cite{rajagopal2024,rezki2024,ordibazar2022}. These streams establish
the relevance of AI for SCRES, but they rarely test learned control
against dense static frontiers on a thesis-grounded resilience metric.

\begin{table}[htbp]
\centering
\caption{Positioning of the present benchmark relative to adjacent literatures.}
\label{tab:literature_gap}
\begin{tabular}{p{0.25\linewidth}ccccc}
\toprule
Stream & Validated DES & SCRES metric & Dense static frontier & Bottleneck action test & Privileged-observation defense \\
\midrule
DES SCRES models & Yes & Yes & Sometimes & Rare & No \\
RL inventory/control & Rare & No & Sometimes & Rare & Rare \\
DES+ML supply-chain workflows & Sometimes & Rare & Rare & Rare & Rare \\
AI-SCRES / MARL reconfiguration & Sometimes & Sometimes & Rare & Strategic & Rare \\
This paper & Yes & Yes & Yes & Yes & Yes \\
\bottomrule
\end{tabular}
\end{table}

\hypertarget{25-gap-action-space-alignment}{%
\subsection{Gap: action-space alignment and bottleneck authority}\label{25-gap-action-space-alignment}}

The gap is therefore conditional rather than categorical. Prior work
shows that RL can be used in supply chains and that DES can represent
SCRES dynamics. What remains under-tested is whether the learned policy's
action space actually reaches the constraint that governs recovery.
Theory-of-Constraints reasoning treats the binding bottleneck as the
lever that determines system throughput \cite{goldratt1984goal}; RL
agent design determines whether that lever is even available to the
policy. This paper connects the two: Track A tests learned control when
the original buffer/shift surface leaves the downstream recovery
bottleneck mostly fixed, while Track B tests a DES-preserving extension
that exposes downstream dispatch authority. The contribution is a
benchmark-centered diagnosis of when AI-enabled SCRES control has a
convertible recovery frontier and when it does not.
